%% ======================================================================
%% 02_context.tex -- 1. El hardware en contexto
%% ======================================================================
\section{El hardware en contexto}
\label{sec:context}

Los sistemas de equilibrio, como el péndulo invertido, la bola y viga, el \textit{ball and plate} y el banco de actitud de multirrotores, son referencias estándar en la enseñanza y la investigación en control porque son inestables en lazo abierto, baratos e intuitivos a la vista \cite{astrom2008,hauser1992,awtar2002}. En HardwareX se han publicado versiones abiertas de estas referencias, por ejemplo el péndulo de varios brazos sobre un carro \cite{kaheman2023} y el péndulo doble de bajo costo con seguimiento por video \cite{myers2020}, junto con robots móviles educativos abiertos como ROBOTONT \cite{raudmae2023} y PlatROB \cite{balbuena2026}, y manipuladores y cabezas educativas como EduSCARA \cite{clark2026} y BEATRIX \cite{al2024}. Estas plataformas estudian un problema de equilibrio a la vez, sobre una base fija.

En robótica de campo el caso relevante es el opuesto: una carga que debe permanecer nivelada viaja sobre un vehículo que acelera, gira y sube pendientes, y el propio vehículo puede llevar una etapa estabilizada activamente. Son ejemplos los estabilizadores de cámara en vehículos terrestres, las bandejas estabilizadas para objetos delicados \cite{selvaggio2023} y los vehículos aéreos que equilibran una carga inestable \cite{hehn2011}. En estos sistemas la perturbación que genera la base se propaga por cada etapa; cada estabilizador atenúa una parte y añade su propia dinámica, retardos y saturaciones. Los proyectos anteriores del curso en nuestra universidad abordaron una sola etapa estabilizada (un cardán de dos ejes sobre un robot diferencial, que detectaba la bola con una pantalla resistiva o con una cámara). Hasta donde sabemos, no existe una plataforma abierta y documentada en la que la propagación a través de \emph{dos} etapas con control activo, una de ellas un vehículo aéreo, pueda medirse con registros sincronizados \emph{y} predecirse con un modelo validado.

\paragraph{Por qué un dron y no un cardán}
La pregunta natural es por qué la etapa intermedia es un octocóptero y no un cardán motorizado de dos ejes, que también nivelaría el plato. La respuesta cuantitativa está en la Sección~\ref{sec:porquedron}. Un cardán solo puede \emph{girar} la carga; el octocóptero, además de nivelarse con la misma ley de control, puede \emph{sostener su propio peso}. Esa fuerza vertical es la única de la plataforma que no pasa por el contacto rueda--suelo, y en la rampa cambia los límites de la base: con la mitad del peso del módulo en los rotores, la pendiente de vuelco estático al subir pasa de \SI{21.5}{\degree} (medida) a \SI{27.7}{\degree} (prevista), al bajar pasa de \SI{30.5}{\degree} (medida) a \SI{38.2}{\degree} (prevista), y el margen de aceleración en una pendiente de \SI{20}{\degree} se multiplica por 3,5. Al mismo tiempo, el dron es el caso más exigente para la etapa intermedia (péndulo invertido con retardo de actuadores y autoridad de par limitada), lo que convierte a la plataforma en un paso previo, medible y seguro, hacia la manipulación aérea en vuelo libre, que es la línea de investigación del grupo.

AeroTray cubre ese vacío con componentes comerciales. Sus aportes son:
\begin{enumerate}[label=(\roman*)]
  \item una arquitectura reproducible de tres niveles (base terrestre, octocóptero sobre una rótula esférica pasiva con guía vertical, \textit{ball and plate}) con interfaces documentadas entre tres controladores independientes (myRIO, ESP32, ESP32-S3) y un supervisor Raspberry Pi, y un sistema de alimentación con dominios aislados y paro de emergencia (Sección~\ref{sec:description});
  \item un modelo cinemático y dinámico de primeros principios de toda la cadena, desde la velocidad de las ruedas hasta la aceleración de la bola, que incluye la base no holonómica en pendiente, el octocóptero como péndulo invertido sobre un pivote acelerado, la asignación de control de ocho rotores y las ecuaciones exactas de una bola que rueda sobre un plato que se traslada y rota; cada ecuación se verificó simbólicamente (Sección~\ref{sec:model} y \ref{app:derivations});
  \item un conjunto de resultados basados en el modelo que explican el comportamiento del prototipo: la identificación del centro de masa del módulo aéreo con telemetría de vuelo estacionario, la justificación cuantitativa del dron en la rampa (vuelco, par y tracción frente a la fracción de peso sostenida), el margen de empuje y la autonomía frente a la tensión de la batería, un mapa de estabilidad del lazo de actitud, una región de atracción por autoridad de par, el límite de estabilidad del lazo de la bola con retardo, la transmisibilidad de las aceleraciones de la base a la bola y los límites de vuelco estático y dinámico de la base (Secciones~\ref{sec:porquedron}--\ref{sec:modelval});
  \item firmware abierto para un controlador de actitud de octocóptero con modo de vuelo en guía, enclavamientos de seguridad y registro a bordo, un controlador de la bola con visión a \SI{187}{fps} en una Raspberry Pi 3B+, y una validación cuantitativa que incluye las iteraciones fallidas, la comparación con el modelo de Isaac Lab usado durante el diseño \cite{mittal2023,schulman2017}, los videos de vuelo sincronizados con la telemetría y todos los datos crudos (Sección~\ref{sec:validation}).
\end{enumerate}
La plataforma está pensada para enseñar integración, modelado y validación en diseño mecatrónico, y como banco de pruebas para controladores en cascada, basados en aprendizaje y con prealimentación de perturbaciones frente a las perturbaciones que induce la base.
